\section{The Problem, by Example}
\label{sec:problem}
Six validators $V_1,\dots,V_6$ confirm payments; at most one of them, say $V_6$, is faulty. A block is confirmed when four validators have given it their \emph{final vote}. Alice signs block $A$ at position 7 of her account, paying Bob, and publishes it.

\para{What a committee change must preserve.}
Suppose the committee is about to change, and the validators agree on a snapshot of all accounts to hand over. Validators $V_1$ to $V_4$ have each signed a final vote for $A$, but the network is slow: no single party has seen all four votes, and in particular the validators building the snapshot have seen at most three. If the snapshot leaves position 7 open, the new committee may confirm a conflicting block $A'$ there. Later the four old votes meet, and $A$ is confirmed too, by a perfectly valid certificate of the old committee. Two conflicting payments are final. Nothing in the agreement on the snapshot went wrong; the snapshot was simply built from a view of the votes that was not the only possible one. We call $A$ \emph{latent-final}: its certificate exists in the released signatures, though nobody has assembled it. The same holds for a certificate that outgoing validators who have not yet stopped may still complete after the snapshot.

The textbook remedy is to stop: end voting, wait until every certificate has either been assembled or become impossible, then snapshot. The wait is as long as the slowest confirmation, at every transition, and during it no payment is confirmed. Section~\ref{sec:eval} measures what this costs: under such a pause, blocks published during a transition wait for it.

RAI's answer is that the snapshot does not need to know whether $A$ is final. It needs to know whether $A$ \emph{could} be final, and that is visible in the reports: a correct validator final-votes only a block it has seen notarized, and says so in its report. Four final voters and five reporters out of six overlap in at least three validators, at least two of them correct, so the snapshot sees $A$ notarized in at least two reports and keeps it. The snapshot does not confirm $A$ (it cannot know), but it \emph{locks} position 7 for $A$, so the new committee cannot confirm $A'$, and $A$'s latent certificate stays compatible with everything that follows.

\para{What a committee change should not cost.}
Nano has no epochs; a representative's weight changes whenever a delegating account moves, read by each node from its own ledger. Two correct nodes can thus see different quorums. Stopping is not an option there either, which is one reason the problem has stayed open. RAI wants both: committee changes with a clear, agreed membership and weight per epoch, and no stop.

\para{Accounts that contradict themselves.}
Now suppose Alice signs two blocks at position 7, $A$ paying Bob and $A'$ paying Carol, and the validators split three against three. No block reaches four. In Nano the position is stuck for good, and so is every later block of Alice's. A deliberate attacker loses access to their own money, which is fine; an honest user with a buggy wallet loses it too, which is not. In RAI the next checkpoint keeps one of the two blocks, under a lock, and every validator that installs it votes only that block at position 7. So the kept payment becomes final within two epochs of normal operation, whatever Alice or a faulty validator does, and Alice's account continues from it. A deliberate attacker gets at most one of their two payments through. The next sections make these statements precise.
